% ---------------------------------------------------------------------------
%  QML Lab research project 02: Trainable Quantum Kernels
%  Build:  latexmk -pdf main.tex   (needs qmlproject.sty, harvardqml.bst, references.bib)
% ---------------------------------------------------------------------------
\documentclass[11pt,a4paper]{article}
\usepackage{qmlproject}
\qmlheader{Project 02: Trainable Quantum Kernels}

\hypersetup{pdftitle={QML Lab project 02: Trainable Quantum Kernels}}

\begin{document}

\begin{center}
  {\small\textcolor{black!55}{QML Lab \textbullet\ research project 02}}\\[0.5em]
  {\LARGE\bfseries Trainable Quantum Kernels}\\[0.6em]
  {\large Learning a feature map whose kernel matches the labels, and measuring what that buys.}\\[0.6em]
  {\small Living document \textbullet\ last updated \today}
\end{center}

\begin{projectcard}
\small
\begin{tabularx}{\linewidth}{@{}>{\bfseries}l X >{\bfseries}l X@{}}
Quadrant  & \CQ\ (classical data, quantum processor) & Track & Kernels \\
Difficulty & Medium & Tools & PennyLane, scikit-learn \\
Builds on & Modules 3 and 6 of the QML Lab & Workspace & \path{qml_projects/trainable-quantum-kernels/} \\
\end{tabularx}
\end{projectcard}

\paragraph{Summary.}
A fixed feature map is a guess about which data points should look similar. This project makes the
feature map trainable and optimises it to maximise kernel--target alignment, a cheap proxy for how
well an SVM will perform. The question is whether alignment training improves test accuracy over a
fixed ZZ map and a tuned RBF kernel, and at what cost in circuit evaluations and overfitting.

% ===========================================================================
\section{Background and key concepts}
% ===========================================================================
\subsection{Quantum embedding kernels}
A quantum embedding kernel uses a parameterised feature map $\ket{\psi_\theta(x)}=U_\theta(x)\ket{0}$,
\begin{equation}
  \kappa_\theta(x,x')=|\braket{\psi_\theta(x)|\psi_\theta(x')}|^2 .
\end{equation}
With fixed $\theta$ this is an ordinary quantum kernel (Module 6). Training $\theta$ adapts the
geometry of the feature space to the task. \citet{hubregtsen2022} studied this on near-term hardware,
including the effect of noise on the kernel matrix.

\subsection{Kernel--target alignment}
Training a full SVM inside an optimisation loop is expensive. Kernel--target alignment
\citep{cristianini2001} measures how close the Gram matrix $K$ is to the ideal kernel $yy^\top$ for
labels $y\in\{-1,+1\}^N$:
\begin{equation}
  A(K,y)=\frac{\langle K,yy^\top\rangle_F}{\lVert K\rVert_F\,\lVert yy^\top\rVert_F}
        =\frac{\sum_{ij}y_iy_jK_{ij}}{N\sqrt{\sum_{ij}K_{ij}^2}} .
\end{equation}
$A=1$ means same-class points are maximally similar and different-class points dissimilar.
Gradients $\partial A/\partial\theta$ follow from parameter-shift gradients of each kernel entry
(Module 3). Each step uses a random batch of $B$ points, so it costs $O(B^2)$ kernel evaluations
rather than $O(N^2)$.

\subsection{Inductive bias}
\citet{kubler2021} show that quantum kernels generalise only if their inductive bias, the functions
they find easy to represent, matches the task; very expressive embeddings spread their capacity
over an exponentially large space and fail to generalise. Training the embedding is one way to
create a useful bias, but it also adds parameters that can overfit.

% ===========================================================================
\section{Research questions}
% ===========================================================================
The project is open-ended: these questions are the starting point, not the finish line. Add new
ones to the README and the journal as they come up.

\begin{rqlist}
  \item Does higher training alignment translate into higher test accuracy, and where does the correlation break down?
  \item How many trainable parameters can the embedding have before the trained kernel overfits?
  \item Can alignment be trained on a small subset (10--20\,\%) of the data and still transfer to the full training set?
\end{rqlist}

% ===========================================================================
\section{Method}
% ===========================================================================
\subsection{Data}
Two moons and concentric circles (scikit-learn) for intuition, then breast cancer (30 features,
PCA to 4) and a two-class subset of Iris. Features are standardised and scaled to $[0,\pi]$.
Every data set gets a fixed train/validation/test split per seed.

\subsection{Models}
A layered embedding on $n=4$ qubits: each layer applies angle encoding of $x$ (features repeated
over qubits if $d<n$), trainable $R_Z$/$R_Y$ rotations and a ring of CNOTs. $\theta$ is trained by
gradient ascent on the alignment of random batches ($B=8$--$16$); the final kernel matrix is then
handed to a precomputed-kernel SVM.

\subsection{Baselines}
The same embedding with random, untrained $\theta$; the fixed ZZ feature map; and an RBF kernel whose
width is also chosen by alignment, so that both sides get the same tuning method.

\subsection{Metrics}
\begin{itemize}
  \item Training and test alignment; test accuracy of the SVM.
  \item Number of kernel evaluations used for training.
  \item Alignment and accuracy as a function of training steps (learning curves).
\end{itemize}

% ===========================================================================
\section{Planned experiments}
% ===========================================================================
Each experiment becomes one \texttt{[[experiments]]} entry in the workspace, so its runs, parameters
and metrics are stored and compared in the portal.

\begin{table}[h]
\centering\small
\begin{tabularx}{\linewidth}{@{}l>{\raggedright}X>{\raggedright}X>{\raggedright\arraybackslash}X@{}}
\toprule
ID & Experiment & What varies & Main metrics \\
\midrule
  E1 & Alignment training on two moons & training steps (0--300) & alignment, test accuracy \\
  E2 & Embedding size & layers $L\in\{1,2,3,4\}$ & train/test gap \\
  E3 & Batch size & $B\in\{4,8,16,32\}$ & alignment noise, cost \\
  E4 & Real data & breast cancer, Iris & test accuracy vs RBF \\
  E5 & Subset transfer & alignment trained on 10\% of data & test accuracy \\
\bottomrule
\end{tabularx}
\caption{Experiment plan. Results go into the journal and the results log (\cref{sec:log}).}
\end{table}

% ===========================================================================
\section{Running it in the QML Lab}
% ===========================================================================
The workspace lives in \path{qml_projects/trainable-quantum-kernels/}. Scripts import three helpers from
\path{qml_lab/plots.py}:
\begin{itemize}
  \item \texttt{params()} reads the parameters from the portal's run form (defaults when run by hand);
  \item \texttt{save()} stores figures with the run;
  \item \texttt{log\_metrics()} records numbers that appear in the run table.
\end{itemize}

\begin{lstlisting}
qml_projects/trainable-quantum-kernels/
  project.toml          metadata + experiments
  README.md             goal, open questions, ideas
  experiments/          runnable scripts
  docs/main.tex         this document
\end{lstlisting}

First experiment: This experiment is \emph{planned}: write the script, add the block below to \texttt{project.toml}, then run \texttt{python manage.py sync\_projects}.

\begin{lstlisting}[language={}]
[[experiments]]
slug = "alignment-moons"
title = "Alignment-trained embedding kernel on two moons"
script = "experiments/e01_alignment_moons.py"
figure = "alignment.png"

[experiments.params]
n_qubits = 4
n_layers = 2
steps = 200
batch = 10
lr = 0.05
seed = 0
\end{lstlisting}

\begin{lstlisting}[language=bash]
# from the portal folder: pick up changes to project.toml
python manage.py sync_projects
cd qml_projects/trainable-quantum-kernels
# run by hand with the default parameters
PYTHONPATH=../../qml_lab python experiments/e01_alignment_moons.py
\end{lstlisting}

% ===========================================================================
\section{Quality bar}
% ===========================================================================
\begin{itemize}
  \item A classical baseline with a similar number of trainable parameters, tuned with the same care \citep{bowles2024}.
  \item At least five random seeds per configuration, reported as mean $\pm$ standard deviation.
  \item How results change with the number of qubits.
  \item At least one run under a realistic noise model.
  \item Choose hyperparameters on a validation split; never tune alignment on the test set.
\end{itemize}

% ===========================================================================
\section{Open directions}
% ===========================================================================
\begin{itemize}
  \item Centred alignment, which is less sensitive to class imbalance.
  \item Multi-class alignment with one-hot label kernels.
  \item Effect of hardware noise on trained kernels, following \citet{hubregtsen2022}.
  \item Compare alignment with the geometric difference of \citet{huang2021power} as a selection criterion.
\end{itemize}

% ===========================================================================
\section{Risks and limitations}
% ===========================================================================
\begin{itemize}
  \item Alignment is a proxy: high alignment on training batches can coexist with poor test accuracy.
  \item Each step costs $O(B^2)$ kernel evaluations, each needing parameter-shift gradients.
  \item Too many embedding parameters overfit small data sets.
\end{itemize}

% ===========================================================================
\section{Results log}
\label{sec:log}
% ===========================================================================
Summarise findings here as they become stable; day-to-day notes belong in the portal journal.

\begin{table}[h]
\centering\small
\begin{tabularx}{\linewidth}{@{}l l >{\raggedright}X >{\raggedright\arraybackslash}X@{}}
\toprule
Date & Run & Finding & Next step \\
\midrule
 & & & \\
 & & & \\
 & & & \\
\bottomrule
\end{tabularx}
\end{table}

\bibliographystyle{harvardqml}
\bibliography{references}

\end{document}
